% Fixed-coordinate diagrams use separate lanes so arrows do not collide.
\hbadness=10000
\vbadness=10000
\hfuzz=2pt
\vfuzz=2pt
\begin{frame}[plain]
\begin{tikzpicture}[remember picture,overlay]
  \node[rounded corners=4pt,fill=blue!65!black,drop shadow={shadow xshift=2pt,shadow yshift=-2pt,opacity=.45},minimum width=15.25cm,minimum height=1.05cm] at ([yshift=-.95cm]current page.north) {};
  \node[text=white,font=\fontsize{25}{28}\selectfont] at ([yshift=-.95cm]current page.north) {PRISMSPACE};
  \fill[prismNavy] (current page.south west) rectangle ([yshift=.45cm]current page.south east);
  \fill[blue!65!black] ([xshift=5.33cm]current page.south west) rectangle ([xshift=10.66cm,yshift=.45cm]current page.south west);
  \node[anchor=center,text=white,font=\scriptsize] at ([xshift=2.66cm,yshift=.22cm]current page.south west) {Group 01 (CEK)};
  \node[anchor=center,text=white,font=\scriptsize] at ([xshift=8cm,yshift=.22cm]current page.south west) {BTech Project};
  \node[anchor=center,text=white,font=\scriptsize] at ([xshift=13.2cm,yshift=.22cm]current page.south west) {September 3, 2026\hspace{1.3cm}1 / \inserttotalframenumber};
  \node[font=\Large\bfseries] at ([yshift=7.20cm]current page.south) {Project Group 01};
  \node[font=\normalsize] at ([yshift=5.90cm]current page.south) {\begin{tabular}{r@{\quad}l}\textbf{NOBIN SIJO} & \textbf{(KNP23CS082)} \\ \textbf{PRANAV P} & \textbf{(KNP23CS085)} \\ \textbf{PRASANTH P} & \textbf{(KNP23CS086)} \\ \textbf{SANKAR NARAYAN S} & \textbf{(KNP23CS094)}\end{tabular}};
  \node[font=\normalsize,align=center] at ([yshift=4.00cm]current page.south) {Under the Guidance of\\[.08cm]Dr. Leena Silvoster M\\ Professor\\College of Engineering, Karunagappally};
  \node at ([yshift=2.25cm]current page.south) {\includegraphics[height=0.95cm]{cek-logo.png}};
  \node[font=\small,align=center] at ([yshift=1.05cm]current page.south) {Dept. of Computer Science \& Engineering\\College of Engineering, Karunagappally};
\end{tikzpicture}
\end{frame}

\begin{frame}{Presentation roadmap}\tableofcontents\end{frame}

\section{Introduction}

\begin{frame}{Introduction}
\vspace{0.25cm}
\begin{block}{Overview}
\normalsize
While modern development tasks span code, browsers, databases, and \textbf{APIs} (Application Programming Interfaces), uncoordinated multi-agent workflows suffer from context fragmentation, runaway loops, and tool hallucinations. To address these challenges, PrismSpace introduces a unified developer workspace featuring \textbf{multi-agent orchestration}, sandboxed \textbf{MCP} (Model Context Protocol) tool execution with strict \textbf{ReAct} (Reasoning and Acting) loops, \textbf{HITL} (Human-in-the-Loop) safety, and \textbf{ORPO} (Odds Ratio Preference Optimization) preference alignment.
\end{block}
\end{frame}

\section{Objectives}

\begin{frame}{Objectives}
\vspace{0.25cm}
\begin{block}{Core Objectives}
\normalsize
The primary objective is to enable structured task execution through \textbf{DAG} (Directed Acyclic Graph)-based task decomposition, \textbf{MCP} (Model Context Protocol)-enabled multi-agent \textbf{ReAct} (Reasoning and Acting) orchestration, and \textbf{ORPO} (Odds Ratio Preference Optimization) preference optimization. Additionally, the system ensures reliable and cost-effective operation by incorporating supervised \textbf{LLM} (Large Language Model) routing alongside Isolation Forest anomaly detection and human-in-the-loop validation.
\end{block}
\end{frame}

\section{Problem statement}

\begin{frame}{Problem Statement}
\vspace{0.25cm}
\begin{block}{Problem Formulation}
\normalsize
Because ungrounded \textbf{LLM} (Large Language Model) agents risk hallucinating tool results and executing unsafe actions, the system takes developer task prompts, workspace states, and \textbf{MCP} (Model Context Protocol)/\textbf{HITL} (Human-in-the-Loop) configurations as its primary \textbf{inputs}. Through an orchestration \textbf{method} integrating sandboxed \textbf{MCP} \textbf{ReAct} (Reasoning and Acting) loops, \textbf{ORPO} (Odds Ratio Preference Optimization), supervised \textbf{LLM} routing, and Isolation Forest anomaly detection with \textbf{HITL} oversight, it converts these tasks into dependable \textbf{outputs} consisting of verified artifacts, \textbf{ReAct} execution traces, and cost telemetry.
\end{block}
\end{frame}

\section{Literature review}

\begin{frame}[t]{Multi-Objective Reward Generalization in Deep Reinforcement Learning}
\vspace{-0.25cm}
\centering
\adjustbox{width=\textwidth, max totalheight=0.785\textheight, keepaspectratio}{%
\small
\renewcommand{\arraystretch}{1.12}
\setlength{\tabcolsep}{4pt}
\begin{tabular}{|p{3.1cm}|p{6.0cm}|p{6.0cm}|}
\hline
\textbf{Citation} & \textbf{Summary} & \textbf{Advantages \& Disadvantages} \\
\hline
\parbox[t]{\linewidth}{\raggedright\small
\textbf{F. Cornalba, C. Disselkamp, D. Scassola, and C. Helf}\newline\vspace{0.2em}
\textit{Neural Computing and Applications}\newline\vspace{0.2em}
\textbf{Year:} 2024\newline\vspace{0.2em}
\textbf{Publisher:} Springer Open Access}
&
Investigates \textbf{Multi-Objective Deep Reinforcement Learning (\textbf{MODRL})} for optimizing competing operational objectives (e.g., return vs. risk). Introduces reward generalization techniques that train pareto-optimal policies capable of dynamically adapting to shifting stakeholder preferences without retraining from scratch.
&
\textbf{Advantages:}
\begin{itemize}\itemsep0.05em\topsep0.1em\topsep0.1em
\item Provides structured Pareto-optimal policy exploration across conflicting objectives.
\item Generalizes across parameterized reward weights without full policy retraining.
\item Directly informs multi-objective routing (latency vs. cost vs. quality) in PrismSpace.
\end{itemize}
\vspace{-0.2em}
\textbf{Disadvantages:}
\begin{itemize}\itemsep0.05em\topsep0.1em\topsep0.1em
\item High initial training sample complexity across multi-dimensional reward spaces.
\item Evaluated primarily on continuous financial time-series data.
\item Sensitive to hyperparameter tuning in policy entropy regularization.
\end{itemize} \\
\hline
\end{tabular}
}
\end{frame}

\begin{frame}[t]{Average Reward Adjusted Discounted Reinforcement Learning}
\vspace{-0.25cm}
\centering
\adjustbox{width=\textwidth, max totalheight=0.785\textheight, keepaspectratio}{%
\small
\renewcommand{\arraystretch}{1.12}
\setlength{\tabcolsep}{4pt}
\begin{tabular}{|p{3.1cm}|p{6.0cm}|p{6.0cm}|}
\hline
\textbf{Citation} & \textbf{Summary} & \textbf{Advantages \& Disadvantages} \\
\hline
\parbox[t]{\linewidth}{\raggedright\small
\textbf{M. Schneckenreither and G. Moser}\newline\vspace{0.2em}
\textit{Neural Computing and Applications}\newline\vspace{0.2em}
\textbf{Year:} 2025\newline\vspace{0.2em}
\textbf{Publisher:} Springer Open Access}
&
Proposes \textbf{Average-Reward Adjusted Discounted \textbf{RL} (Reinforcement Learning) (ARAD-\textbf{RL})}, bridging the gap between discounted and average reward criteria in infinite-horizon \textbf{MDPs} (Markov Decision Processes). Addresses systematic value estimation bias in standard discounted Q-learning by dynamically subtracting the estimated gain rate during Bellman updates.
&
\textbf{Advantages:}
\begin{itemize}\itemsep0.05em\topsep0.1em
\item Eliminates infinite-horizon return explosion while preserving policy convergence.
\item Improves long-term steady-state throughput in continuous scheduling tasks.
\item Theoretically sound formulation for continuous agent lifecycle optimization.
\end{itemize}
\vspace{-0.2em}
\textbf{Disadvantages:}
\begin{itemize}\itemsep0.05em\topsep0.1em
\item Requires accurate empirical estimation of average reward baseline.
\item Slow convergence in highly non-stationary transition environments.
\item Adds variance to early-stage \textbf{TD} (Temporal Difference) error gradient estimates.
\end{itemize} \\
\hline
\end{tabular}
}
\end{frame}

\begin{frame}[t]{Deep Reinforcement Learning for Multi-Agent Coordination}
\vspace{-0.25cm}
\centering
\adjustbox{width=\textwidth, max totalheight=0.785\textheight, keepaspectratio}{%
\small
\renewcommand{\arraystretch}{1.12}
\setlength{\tabcolsep}{4pt}
\begin{tabular}{|p{3.1cm}|p{6.0cm}|p{6.0cm}|}
\hline
\textbf{Citation} & \textbf{Summary} & \textbf{Advantages \& Disadvantages} \\
\hline
\parbox[t]{\linewidth}{\raggedright\small
\textbf{K. Aina and S. Ha}\newline\vspace{0.2em}
\textit{Artificial Life and Robotics}\newline\vspace{0.2em}
\textbf{Year:} 2026\newline\vspace{0.2em}
\textbf{Publisher:} Springer Open Access}
&
Develops decentralized \textbf{MARL} (Multi-Agent Reinforcement Learning) algorithms for multi-agent coordination in constrained environments. Employs \textbf{spatial-temporal observation encoding} and localized interaction graphs to coordinate joint movements and resolve deadlocks without requiring centralized global state knowledge.
&
\textbf{Advantages:}
\begin{itemize}\itemsep0.05em\topsep0.1em
\item Resolves shared-resource contention and deadlocks via decentralized policies.
\item Spatial-temporal graph attention scales linearly with agent density.
\item Applicable to concurrent tool-use contention across worker agents.
\end{itemize}
\vspace{-0.2em}
\textbf{Disadvantages:}
\begin{itemize}\itemsep0.05em\topsep0.1em
\item Tested in physical robotics grids rather than software \textbf{API} (Application Programming Interface) environments.
\item Vulnerable to communication packet dropouts and sensor latency.
\item Does not account for semantic language-based coordination.
\end{itemize} \\
\hline
\end{tabular}
}
\end{frame}

\begin{frame}[t]{Open-Ended Coordination Using Modular Open Policies}
\vspace{-0.25cm}
\centering
\adjustbox{width=\textwidth, max totalheight=0.785\textheight, keepaspectratio}{%
\small
\renewcommand{\arraystretch}{1.12}
\setlength{\tabcolsep}{4pt}
\begin{tabular}{|p{3.1cm}|p{6.0cm}|p{6.0cm}|}
\hline
\textbf{Citation} & \textbf{Summary} & \textbf{Advantages \& Disadvantages} \\
\hline
\parbox[t]{\linewidth}{\raggedright\small
\textbf{D. Rother, J. Pajarinen, J. Peters, and T. H. Weisswange}\newline\vspace{0.2em}
\textit{Autonomous Agents and Multi-Agent Systems}\newline\vspace{0.2em}
\textbf{Year:} 2025\newline\vspace{0.2em}
\textbf{Publisher:} Springer Open Access}
&
Presents a framework for \textbf{open-ended ad hoc multi-agent coordination}. Utilizes \textbf{Modular Open Policies (\textbf{MOP})} where dynamic agent teams dynamically infer unseen teammate skill representations and synthesize blended coordination strategies without requiring pre-coordinated joint training.
&
\textbf{Advantages:}
\begin{itemize}\itemsep0.05em\topsep0.1em
\item Enables zero-shot ad hoc coordination with unfamiliar teammate agent policies.
\item Modular policy decomposition allows plug-and-play agent role assignment.
\item Directly adopted in PrismSpace for dynamic worker swarm scaling.
\end{itemize}
\vspace{-0.2em}
\textbf{Disadvantages:}
\begin{itemize}\itemsep0.05em\topsep0.1em
\item Requires comprehensive pre-trained skill policy libraries for base roles.
\item Evaluated primarily on benchmark games (Overcooked, \textbf{LBF} (Level-Based Foraging)).
\item Action blending may introduce transient sub-optimal execution steps.
\end{itemize} \\
\hline
\end{tabular}
}
\end{frame}

\begin{frame}[t]{Improving Scalability of \textbf{MADRL} (Multi-Agent Deep Reinforcement Learning) with Suboptimal Human Knowledge}
\vspace{-0.08cm}
\centering
\adjustbox{width=\textwidth, max totalheight=0.785\textheight, keepaspectratio}{%
\small
\renewcommand{\arraystretch}{1.12}
\setlength{\tabcolsep}{4pt}
\begin{tabular}{|p{3.1cm}|p{6.0cm}|p{6.0cm}|}
\hline
\textbf{Citation} & \textbf{Summary} & \textbf{Advantages \& Disadvantages} \\
\hline
\parbox[t]{\linewidth}{\raggedright\small
\textbf{D. Liu, F. Ren, J. Yan, G. Su, W. Gu, and S. Kato}\newline\vspace{0.2em}
\textit{Autonomous Agents and Multi-Agent Systems}\newline\vspace{0.2em}
\textbf{Year:} 2026\newline\vspace{0.2em}
\textbf{Publisher:} Springer Open Access}
&
Addresses sample inefficiency in large-scale \textbf{Multi-Agent DRL} (Deep Reinforcement Learning) by integrating heuristic human guidance into exploration. Proposes an adaptive filtering mechanism that extracts actionable macro-actions from noisy or suboptimal human demonstration traces, accelerating multi-agent policy learning.
&
\textbf{Advantages:}
\begin{itemize}\itemsep0.05em\topsep0.1em
\item Accelerates policy exploration by up to 4x using noisy demonstration priors.
\item Robust to suboptimal or conflicting human heuristic rules.
\item Informs PrismSpace's bootstrap initialization from logged developer sessions.
\end{itemize}
\vspace{-0.2em}
\textbf{Disadvantages:}
\begin{itemize}\itemsep0.05em\topsep0.1em
\item Quality of learned policies depends on initial heuristic coverage.
\item Requires complex threshold tuning to transition away from demonstrations.
\item High memory footprint for maintaining dual replay buffers.
\end{itemize} \\
\hline
\end{tabular}
}
\end{frame}

\begin{frame}[t]{$\varepsilon$-Retraining Reinforcement Learning Algorithms}
\vspace{-0.25cm}
\centering
\adjustbox{width=\textwidth, max totalheight=0.785\textheight, keepaspectratio}{%
\small
\renewcommand{\arraystretch}{1.12}
\setlength{\tabcolsep}{4pt}
\begin{tabular}{|p{3.1cm}|p{6.0cm}|p{6.0cm}|}
\hline
\textbf{Citation} & \textbf{Summary} & \textbf{Advantages \& Disadvantages} \\
\hline
\parbox[t]{\linewidth}{\raggedright\small
\textbf{L. Marzari, C. Liu, P. L. Donti, and E. Marchesini}\newline\vspace{0.2em}
\textit{Autonomous Agents and Multi-Agent Systems}\newline\vspace{0.2em}
\textbf{Year:} 2026\newline\vspace{0.2em}
\textbf{Publisher:} Springer Open Access}
&
Investigates bounded policy fine-tuning through \textbf{$\varepsilon$-retraining}, preventing catastrophic forgetting when deep \textbf{RL} agents adapt to non-stationary environments. Establishes formal Lipschitz bounds on policy drift, ensuring safe continual learning in deployed agent pipelines.
&
\textbf{Advantages:}
\begin{itemize}\itemsep0.05em\topsep0.1em
\item Provides mathematical guarantees against catastrophic policy divergence.
\item Enables continuous telemetry-driven model fine-tuning in production.
\item Crucial for stability in PrismSpace's ongoing user interaction learning.
\end{itemize}
\vspace{-0.2em}
\textbf{Disadvantages:}
\begin{itemize}\itemsep0.05em\topsep0.1em
\item Constrained optimization increases computation per gradient step.
\item Conservative update bounds may slow adaptation to drastic distribution shifts.
\item Theoretical guarantees rely on local smoothness assumptions.
\end{itemize} \\
\hline
\end{tabular}
}
\end{frame}

\begin{frame}[t]{Deep Multiagent Reinforcement Learning: Challenges and Directions}
\vspace{-0.25cm}
\centering
\adjustbox{width=\textwidth, max totalheight=0.785\textheight, keepaspectratio}{%
\small
\renewcommand{\arraystretch}{1.12}
\setlength{\tabcolsep}{4pt}
\begin{tabular}{|p{3.1cm}|p{6.0cm}|p{6.0cm}|}
\hline
\textbf{Citation} & \textbf{Summary} & \textbf{Advantages \& Disadvantages} \\
\hline
\parbox[t]{\linewidth}{\raggedright\small
\textbf{A. Wong, T. B\"ack, A. V. Kononova, and A. Plaat}\newline\vspace{0.2em}
\textit{Artificial Intelligence Review}\newline\vspace{0.2em}
\textbf{Year:} 2023\newline\vspace{0.2em}
\textbf{Publisher:} Springer Open Access}
&
Surveys \textbf{Deep Multiagent Reinforcement Learning (\textbf{MARL})}, formalizing foundational challenges including non-stationarity, multi-agent credit assignment, and combinatorial action explosion. Evaluates \textbf{Centralized Training with Decentralized Execution (\textbf{CTDE})} and communication architectures for modular orchestration.
&
\textbf{Advantages:}
\begin{itemize}\itemsep0.05em\topsep0.1em
\item Comprehensive taxonomy of credit assignment and \textbf{CTDE} paradigms.
\item Establishes formal theoretical criteria for communication protocols.
\item Provides the theoretical foundation for PrismSpace's modular agent architecture.
\end{itemize}
\vspace{-0.2em}
\textbf{Disadvantages:}
\begin{itemize}\itemsep0.05em\topsep0.1em
\item Survey paper without empirical experimental evaluation.
\item Focuses on discrete game simulations rather than \textbf{API}-driven software tasks.
\item Does not cover \textbf{LLM} reasoning agents or human-in-the-loop approval gates.
\end{itemize} \\
\hline
\end{tabular}
}
\end{frame}

\begin{frame}[t]{\textbf{HIAT} (Human-in-the-Loop with Auxiliary Tasks): Human-in-the-Loop \textbf{RL} with Auxiliary Tasks}
\vspace{-0.08cm}
\centering
\adjustbox{width=\textwidth, max totalheight=0.785\textheight, keepaspectratio}{%
\small
\renewcommand{\arraystretch}{1.12}
\setlength{\tabcolsep}{4pt}
\begin{tabular}{|p{3.1cm}|p{6.0cm}|p{6.0cm}|}
\hline
\textbf{Citation} & \textbf{Summary} & \textbf{Advantages \& Disadvantages} \\
\hline
\parbox[t]{\linewidth}{\raggedright\small
\textbf{B. Niu, B. Luo, Y. Cui, X. Xu, Y. Zhao, and Y. Feng}\newline\vspace{0.2em}
\textit{Artificial Intelligence Review}\newline\vspace{0.2em}
\textbf{Year:} 2025\newline\vspace{0.2em}
\textbf{Publisher:} Springer Open Access}
&
Introduces \textbf{HIAT (Human-in-the-Loop with Auxiliary Tasks)}, a framework leveraging human intervention signals as auxiliary self-supervised learning objectives. Models human corrections as feedback gradients that shape internal feature representations, minimizing human intervention frequency over time.
&
\textbf{Advantages:}
\begin{itemize}\itemsep0.05em\topsep0.1em
\item Maximizes learning efficiency from scarce human intervention signals.
\item Significantly reduces human cognitive load through auxiliary reward shaping.
\item Informs PrismSpace's approval gate telemetry and learning from user edits.
\end{itemize}
\vspace{-0.2em}
\textbf{Disadvantages:}
\begin{itemize}\itemsep0.05em\topsep0.1em
\item Requires consistent human feedback to avoid biased auxiliary representations.
\item Intervention timing latency can cause temporal mis-alignment in \textbf{RL} updates.
\item Higher computational overhead during auxiliary task representation training.
\end{itemize} \\
\hline
\end{tabular}
}
\end{frame}

\begin{frame}[t]{\textbf{ORPO} (Odds Ratio Preference Optimization): Monolithic Preference Optimization without Reference Model}
\vspace{-0.22cm}
\centering
\adjustbox{width=\textwidth, max totalheight=0.785\textheight, keepaspectratio}{%
\small
\renewcommand{\arraystretch}{1.12}
\setlength{\tabcolsep}{4pt}
\begin{tabular}{|p{3.1cm}|p{6.0cm}|p{6.0cm}|}
\hline
\textbf{Citation} & \textbf{Summary} & \textbf{Advantages \& Disadvantages} \\
\hline
\parbox[t]{\linewidth}{\raggedright\small
\textbf{J. Hong, N. Lee, and J. Thorne}\newline\vspace{0.2em}
\textit{\textbf{EMNLP} / arXiv:2403.07691}\newline\vspace{0.2em}
\textbf{Year:} 2024\newline\vspace{0.2em}
\textbf{Area:} \textbf{LLM} Preference Alignment}
&
\textbf{Core Idea (Single-Phase Alignment):}\newline
Instead of multi-step \textbf{RLHF}, \textbf{ORPO} teaches desired answers while penalizing bad answers in one unified equation:
\vspace{-0.15cm}
\[ \mathcal{L}_{\text{ORPO}} = \mathcal{L}_{\text{SFT}} + \lambda \cdot \mathcal{L}_{\text{OR}} \]
\vspace{-0.30cm}
\[ \mathcal{L}_{\text{OR}} = -\log \sigma \left( \log \frac{\text{odds}(y_w|x)}{\text{odds}(y_l|x)} \right) \]
\vspace{-0.20cm}
\textbf{Key Terms:}
\begin{itemize}\itemsep0.02em\topsep0.05em
  \item $\mathcal{L}_{\text{SFT}}$: Base supervised fine-tuning loss (correct code/syntax).
  \item $\mathcal{L}_{\text{OR}}, \lambda$: Odds ratio penalty term and penalty weight.
  \item $x, y_w, y_l$: Prompt, winning action path, and losing action path.
\end{itemize}
&
\textbf{Advantages:}
\begin{itemize}\itemsep0.05em\topsep0.1em
\item Reference-model-free: eliminates maintaining frozen $\pi_{\text{ref}}$ in memory.
\item Monolithic single-phase training avoids unstable \textbf{PPO} policy optimization.
\item Directly aligns PrismSpace agent trajectories on paired feedback (\textit{chosen} vs. \textit{rejected}).
\end{itemize}
\vspace{-0.2em}
\textbf{Disadvantages:}
\begin{itemize}\itemsep0.05em\topsep0.1em
\item Requires high-quality pairwise preference datasets ($y_w, y_l$).
\item Sensitive to $\lambda$ hyperparameter balance.
\end{itemize} \\
\hline
\end{tabular}
}
\end{frame}

\begin{frame}[t]{Reinforcement Online Learning to Rank with Unbiased Reward Shaping}
\vspace{-0.25cm}
\centering
\adjustbox{width=\textwidth, max totalheight=0.785\textheight, keepaspectratio}{%
\small
\renewcommand{\arraystretch}{1.12}
\setlength{\tabcolsep}{4pt}
\begin{tabular}{|p{3.1cm}|p{6.0cm}|p{6.0cm}|}
\hline
\textbf{Citation} & \textbf{Summary} & \textbf{Advantages \& Disadvantages} \\
\hline
\parbox[t]{\linewidth}{\raggedright\small
\textbf{S. Zhuang, Z. Qiao, and G. Zuccon}\newline\vspace{0.2em}
\textit{Information Retrieval Journal}\newline\vspace{0.2em}
\textbf{Year:} 2023\newline\vspace{0.2em}
\textbf{Publisher:} Springer Open Access}
&
Develops an \textbf{Online Learning to Rank (\textbf{OLTR})} framework using counterfactual reinforcement learning. Employs \textbf{inverse propensity weighting (\textbf{IPW})} and unbiased reward shaping to correct for position, presentation, and click bias in user interaction logs during continuous ranker optimization.
&
\textbf{Advantages:}
\begin{itemize}\itemsep0.05em\topsep0.1em
\item Guarantees unbiased gradient estimation from observational implicit feedback.
\item Operates entirely online without requiring manual document relevance labels.
\item Directly informs PrismSpace's context retrieval ranking feedback loop.
\end{itemize}
\vspace{-0.2em}
\textbf{Disadvantages:}
\begin{itemize}\itemsep0.05em\topsep0.1em
\item High propensity score variance can destabilize gradient updates.
\item Requires large volumes of query-click logs to converge.
\item Assumes position-based user click models that may not hold in code search.
\end{itemize} \\
\hline
\end{tabular}
}
\end{frame}

\begin{frame}[t]{Knowledge Learning for Securing Workflow Scheduling in \textbf{MEC} (Mobile Edge Computing)}
\vspace{-0.08cm}
\centering
\adjustbox{width=\textwidth, max totalheight=0.785\textheight, keepaspectratio}{%
\small
\renewcommand{\arraystretch}{1.12}
\setlength{\tabcolsep}{4pt}
\begin{tabular}{|p{3.1cm}|p{6.0cm}|p{6.0cm}|}
\hline
\textbf{Citation} & \textbf{Summary} & \textbf{Advantages \& Disadvantages} \\
\hline
\parbox[t]{\linewidth}{\raggedright\small
\textbf{T. M. Ghazal, A. E. Awouda, M. K. Hasan, A. H. A. Rahman, S. Islam, R. A. Saeed, H. Elshafie, A. Iqbal, and A. H. Hussein}\newline\vspace{0.2em}
\textit{Mobile Networks and Applications}\newline\vspace{0.2em}
\textbf{Year:} 2025\newline\vspace{0.2em}
\textbf{Publisher:} Springer Open Access}
&
Proposes a \textbf{Knowledge-Guided Reinforcement Learning (\textbf{KG}-\textbf{RL})} scheduler for complex workflow \textbf{DAGs} in distributed edge-cloud environments. Integrates security constraint validation directly into the state transition model, preventing insecure task placement while optimizing makespan.
&
\textbf{Advantages:}
\begin{itemize}\itemsep0.05em\topsep0.1em
\item Enforces strict security boundary constraints during automated task scheduling.
\item Effectively handles multi-dependency Directed Acyclic Graph (\textbf{DAG}) workflows.
\item Informs PrismSpace's sandboxed tool execution and dependency resolution.
\end{itemize}
\vspace{-0.2em}
\textbf{Disadvantages:}
\begin{itemize}\itemsep0.05em\topsep0.1em
\item Constraint checking introduces scheduling latency for massive \textbf{DAGs}.
\item Assumes static security classification matrices across network nodes.
\item Limited evaluation on non-deterministic tool failure scenarios.
\end{itemize} \\
\hline
\end{tabular}
}
\end{frame}

\begin{frame}[t]{Energy-Aware Cooperative Multi-Fitness Evolutionary Workflow Scheduling}
\vspace{-0.25cm}
\centering
\adjustbox{width=\textwidth, max totalheight=0.785\textheight, keepaspectratio}{%
\small
\renewcommand{\arraystretch}{1.12}
\setlength{\tabcolsep}{4pt}
\begin{tabular}{|p{3.1cm}|p{6.0cm}|p{6.0cm}|}
\hline
\textbf{Citation} & \textbf{Summary} & \textbf{Advantages \& Disadvantages} \\
\hline
\parbox[t]{\linewidth}{\raggedright\small
\textbf{P. Barredo and J. Puente}\newline\vspace{0.2em}
\textit{Natural Computing}\newline\vspace{0.2em}
\textbf{Year:} 2025\newline\vspace{0.2em}
\textbf{Publisher:} Springer Open Access}
&
Presents a \textbf{Cooperative Multi-Fitness Evolutionary Algorithm (\textbf{CMFEA})} for multi-objective workflow scheduling. Optimizes trade-offs between execution time, monetary compute cost, and energy consumption using island-model cooperative coevolution across heterogeneous worker pools.
&
\textbf{Advantages:}
\begin{itemize}\itemsep0.05em\topsep0.1em
\item Pareto-optimal trade-offs across makespan, financial cost, and energy.
\item Cooperative sub-population decomposition prevents premature local convergence.
\item Provides heuristics for PrismSpace's multi-agent model provider selection.
\end{itemize}
\vspace{-0.2em}
\textbf{Disadvantages:}
\begin{itemize}\itemsep0.05em\topsep0.1em
\item Evolutionary search has higher runtime latency than single-pass ML heuristics.
\item Cannot adapt instantaneously to real-time \textbf{API} rate limit drops.
\item Requires parameter tuning across mutation, crossover, and migration rates.
\end{itemize} \\
\hline
\end{tabular}
}
\end{frame}

\begin{frame}[t]{Optimized Task Scheduling Using Hybrid \textbf{DL} (Deep Learning) and Metaheuristics}
\vspace{-0.08cm}
\centering
\adjustbox{width=\textwidth, max totalheight=0.785\textheight, keepaspectratio}{%
\small
\renewcommand{\arraystretch}{1.12}
\setlength{\tabcolsep}{4pt}
\begin{tabular}{|p{3.1cm}|p{6.0cm}|p{6.0cm}|}
\hline
\textbf{Citation} & \textbf{Summary} & \textbf{Advantages \& Disadvantages} \\
\hline
\parbox[t]{\linewidth}{\raggedright\small
\textbf{A. Heirati, M. Fartash, and M. Khalily-Dermany}\newline\vspace{0.2em}
\textit{Neural Processing Letters}\newline\vspace{0.2em}
\textbf{Year:} 2026\newline\vspace{0.2em}
\textbf{Publisher:} Springer Open Access}
&
Combines \textbf{Deep Learning (\textbf{LSTM} (Long Short-Term Memory))} workload prediction with metaheuristic optimization (Harris Hawks Optimization) for proactive workflow dispatching. Predicts task execution durations and memory bottlenecks in advance to prevent node overloading and minimize queue waiting times.
&
\textbf{Advantages:}
\begin{itemize}\itemsep0.05em\topsep0.1em
\item Proactive workload forecasting prevents sudden compute node thrashing.
\item Hybrid pipeline balances search exploration with predictive dispatching.
\item Relevant for predicting token latency and cost before agent invocation.
\end{itemize}
\vspace{-0.2em}
\textbf{Disadvantages:}
\begin{itemize}\itemsep0.05em\topsep0.1em
\item Two-stage pipeline introduces inference overhead before scheduling decisions.
\item \textbf{LSTM} prediction accuracy degrades under sudden irregular workload spikes.
\item Complex offline model retraining requirements.
\end{itemize} \\
\hline
\end{tabular}
}
\end{frame}

\begin{frame}[t]{Cost-Aware Workflow Offloading in Edge-Cloud Computing via \textbf{GA} (Genetic Algorithm)}
\vspace{-0.08cm}
\centering
\adjustbox{width=\textwidth, max totalheight=0.785\textheight, keepaspectratio}{%
\small
\renewcommand{\arraystretch}{1.12}
\setlength{\tabcolsep}{4pt}
\begin{tabular}{|p{3.1cm}|p{6.0cm}|p{6.0cm}|}
\hline
\textbf{Citation} & \textbf{Summary} & \textbf{Advantages \& Disadvantages} \\
\hline
\parbox[t]{\linewidth}{\raggedright\small
\textbf{S. Abdi, M. Ashjaei, and S. Mubeen}\newline\vspace{0.2em}
\textit{The Journal of Supercomputing}\newline\vspace{0.2em}
\textbf{Year:} 2024\newline\vspace{0.2em}
\textbf{Publisher:} Springer Open Access}
&
Proposes a \textbf{Genetic Algorithm-based workflow offloading} framework considering heterogeneous network bandwidth, deadline constraints, and monetary costs. Implements chromosome repair operators to preserve \textbf{DAG} precedence validity throughout mutation and crossover.
&
\textbf{Advantages:}
\begin{itemize}\itemsep0.05em\topsep0.1em
\item Guarantees valid \textbf{DAG} dependency structures via specialized repair operators.
\item Explicitly accounts for network transfer latency between pipeline stages.
\item Informs subtask payload transfer between worker agents in PrismSpace.
\end{itemize}
\vspace{-0.2em}
\textbf{Disadvantages:}
\begin{itemize}\itemsep0.05em\topsep0.1em
\item Computationally heavy for interactive real-time user query response.
\item Assumes deterministic network transfer bandwidth in cloud environments.
\item Does not model dynamic agent failures or retries.
\end{itemize} \\
\hline
\end{tabular}
}
\end{frame}

\begin{frame}[t]{Dynamic Role Discovery and Assignment in Multi-Agent Task Decomposition}
\vspace{-0.25cm}
\centering
\adjustbox{width=\textwidth, max totalheight=0.785\textheight, keepaspectratio}{%
\small
\renewcommand{\arraystretch}{1.12}
\setlength{\tabcolsep}{4pt}
\begin{tabular}{|p{3.1cm}|p{6.0cm}|p{6.0cm}|}
\hline
\textbf{Citation} & \textbf{Summary} & \textbf{Advantages \& Disadvantages} \\
\hline
\parbox[t]{\linewidth}{\raggedright\small
\textbf{Y. Xia, J. Zhu, and L. Zhu}\newline\vspace{0.2em}
\textit{Complex \& Intelligent Systems}\newline\vspace{0.2em}
\textbf{Year:} 2023\newline\vspace{0.2em}
\textbf{Publisher:} Springer Open Access}
&
Proposes \textbf{dynamic role discovery and assignment} that decomposes complex high-level objectives into lower-level subtasks. Uses learned persona embeddings and mutual information to adaptively allocate subtasks according to situational agent capabilities, eliminating rigid static assignments.
&
\textbf{Advantages:}
\begin{itemize}\itemsep0.05em\topsep0.1em
\item Enables flexible subtask allocation without hardcoded role hierarchies.
\item Hierarchical decomposition prevents single-agent cognitive context overload.
\item Serves as core theoretical basis for PrismSpace's dynamic agent routing.
\end{itemize}
\vspace{-0.2em}
\textbf{Disadvantages:}
\begin{itemize}\itemsep0.05em\topsep0.1em
\item Role discovery clustering introduces upfront computational overhead.
\item Performance depends on representation vector accuracy for subtasks.
\item Evaluated primarily on synthetic collaborative coordination benchmarks.
\end{itemize} \\
\hline
\end{tabular}
}
\end{frame}

\begin{frame}[t]{Decision Variable Contribution Based Adaptive Multi-Objective Scheduling}
\vspace{-0.25cm}
\centering
\adjustbox{width=\textwidth, max totalheight=0.785\textheight, keepaspectratio}{%
\small
\renewcommand{\arraystretch}{1.12}
\setlength{\tabcolsep}{4pt}
\begin{tabular}{|p{3.1cm}|p{6.0cm}|p{6.0cm}|}
\hline
\textbf{Citation} & \textbf{Summary} & \textbf{Advantages \& Disadvantages} \\
\hline
\parbox[t]{\linewidth}{\raggedright\small
\textbf{J. Li, L. Xing, W. Zhong, Z. Cai, and F. Hou}\newline\vspace{0.2em}
\textit{Complex \& Intelligent Systems}\newline\vspace{0.2em}
\textbf{Year:} 2023\newline\vspace{0.2em}
\textbf{Publisher:} Springer Open Access}
&
Introduces a \textbf{Decision Variable Contribution (DVC)} adaptive mechanism for evolutionary multi-objective workflow scheduling. Evaluates the differential sensitivity of individual decision variables to focus evolutionary search operators on high-impact scheduling bottlenecks.
&
\textbf{Advantages:}
\begin{itemize}\itemsep0.05em\topsep0.1em
\item Accelerates Pareto-front convergence by prioritizing high-variance variables.
\item Dynamically adapts search intensity across workflow bottleneck stages.
\item Useful for pinpointing high-latency critical path tools in agent workflows.
\end{itemize}
\vspace{-0.2em}
\textbf{Disadvantages:}
\begin{itemize}\itemsep0.05em\topsep0.1em
\item Sensitivity calculation adds per-generation matrix computation cost.
\item Can over-concentrate on dominant objectives at the expense of diversity.
\item Requires calibrated normalization across widely different metric units.
\end{itemize} \\
\hline
\end{tabular}
}
\end{frame}

\begin{frame}[t]{GHQ: Grouped Hybrid Q-Learning for Heterogeneous Multi-Agent RL}
\vspace{-0.25cm}
\centering
\adjustbox{width=\textwidth, max totalheight=0.785\textheight, keepaspectratio}{%
\small
\renewcommand{\arraystretch}{1.12}
\setlength{\tabcolsep}{4pt}
\begin{tabular}{|p{3.1cm}|p{6.0cm}|p{6.0cm}|}
\hline
\textbf{Citation} & \textbf{Summary} & \textbf{Advantages \& Disadvantages} \\
\hline
\parbox[t]{\linewidth}{\raggedright\small
\textbf{X. Yu, Y. Lin, X. Wang, S. Han, and K. Lv}\newline\vspace{0.2em}
\textit{Complex \& Intelligent Systems}\newline\vspace{0.2em}
\textbf{Year:} 2024\newline\vspace{0.2em}
\textbf{Publisher:} Springer Open Access}
&
Develops \textbf{Grouped Hybrid Q-learning (GHQ)} for cooperative heterogeneous multi-agent systems. Clusters agents with similar capabilities into functional cohorts that share value approximations while maintaining decoupled action execution heads for role specialization.
&
\textbf{Advantages:}
\begin{itemize}\itemsep0.05em\topsep0.1em
\item Scales value estimation across large heterogeneous swarms via cohort grouping.
\item Reduces joint action space dimensionality from exponential to block-linear.
\item Directly applicable to categorizing \textbf{LLM} worker agent skill specializations.
\end{itemize}
\vspace{-0.2em}
\textbf{Disadvantages:}
\begin{itemize}\itemsep0.05em\topsep0.1em
\item Cohort clustering boundaries can be rigid if agent capabilities evolve.
\item Intra-group credit assignment remains non-trivial in dense interactions.
\item Requires domain knowledge to establish initial grouping heuristics.
\end{itemize} \\
\hline
\end{tabular}
}
\end{frame}

\begin{frame}[t]{Decentralized Multi-Agent Cooperation via Adaptive Partner Modeling}
\vspace{-0.25cm}
\centering
\adjustbox{width=\textwidth, max totalheight=0.785\textheight, keepaspectratio}{%
\small
\renewcommand{\arraystretch}{1.12}
\setlength{\tabcolsep}{4pt}
\begin{tabular}{|p{3.1cm}|p{6.0cm}|p{6.0cm}|}
\hline
\textbf{Citation} & \textbf{Summary} & \textbf{Advantages \& Disadvantages} \\
\hline
\parbox[t]{\linewidth}{\raggedright\small
\textbf{C. Xu, J. Wang, X. Zhu, Y. Yue, W. Zhou, Z. Liang, and D. Wojtczak}\newline\vspace{0.2em}
\textit{Complex \& Intelligent Systems}\newline\vspace{0.2em}
\textbf{Year:} 2024\newline\vspace{0.2em}
\textbf{Publisher:} Springer Open Access}
&
Tackles environmental non-stationarity in decentralized multi-agent cooperation using \textbf{Adaptive Partner Modeling (APM)}. Instead of relying on expensive continuous inter-agent broadcast communication, agents infer partner policies and intended actions from local observation histories.
&
\textbf{Advantages:}
\begin{itemize}\itemsep0.05em\topsep0.1em
\item Minimizes inter-agent communication bandwidth to near zero during execution.
\item Stabilizes decentralized joint execution despite teammate policy shifts.
\item Directly adopted in PrismSpace for zero-communication agent coordination.
\end{itemize}
\vspace{-0.2em}
\textbf{Disadvantages:}
\begin{itemize}\itemsep0.05em\topsep0.1em
\item Partner inference error propagates if historical observations are noisy.
\item Inference networks introduce local computational and memory overhead.
\item Assumes partners follow relatively stationary parameterized policy spaces.
\end{itemize} \\
\hline
\end{tabular}
}
\end{frame}

\begin{frame}[t]{A Responsible AI Framework: Pipeline Contextualisation}
\vspace{-0.25cm}
\centering
\adjustbox{width=\textwidth, max totalheight=0.785\textheight, keepaspectratio}{%
\small
\renewcommand{\arraystretch}{1.12}
\setlength{\tabcolsep}{4pt}
\begin{tabular}{|p{3.1cm}|p{6.0cm}|p{6.0cm}|}
\hline
\textbf{Citation} & \textbf{Summary} & \textbf{Advantages \& Disadvantages} \\
\hline
\parbox[t]{\linewidth}{\raggedright\small
\textbf{E. Vyhmeister, G. G. Casta\~n\'e, P.-O. \"Ostberg, and S. Thevenin}\newline\vspace{0.2em}
\textit{AI and Ethics}\newline\vspace{0.2em}
\textbf{Year:} 2023\newline\vspace{0.2em}
\textbf{Publisher:} Springer Open Access}
&
Formulates an operational \textbf{Responsible AI Framework} providing pipeline contextualisation across the full model lifecycle. Specifies auditable checkpoints, data lineage tracking, and bias mitigation protocols for autonomous AI decision-making systems.
&
\textbf{Advantages:}
\begin{itemize}\itemsep0.05em\topsep0.1em
\item Holistic governance framework linking technical architecture to ethical safety.
\item Defines concrete telemetry points for data provenance and action auditing.
\item Guides PrismSpace's execution audit logging and human inspection UI.
\end{itemize}
\vspace{-0.2em}
\textbf{Disadvantages:}
\begin{itemize}\itemsep0.05em\topsep0.1em
\item Qualitative framework that requires custom translation into concrete code gates.
\item Audit logging introduces storage overhead and slight execution latency.
\item Does not provide automated formal verification of \textbf{LLM} reasoning traces.
\end{itemize} \\
\hline
\end{tabular}
}
\end{frame}

\begin{frame}[t]{Meaningful Human Control: Actionable Properties for AI Development}
\vspace{-0.25cm}
\centering
\adjustbox{width=\textwidth, max totalheight=0.785\textheight, keepaspectratio}{%
\small
\renewcommand{\arraystretch}{1.12}
\setlength{\tabcolsep}{4pt}
\begin{tabular}{|p{3.1cm}|p{6.0cm}|p{6.0cm}|}
\hline
\textbf{Citation} & \textbf{Summary} & \textbf{Advantages \& Disadvantages} \\
\hline
\parbox[t]{\linewidth}{\raggedright\small
\textbf{L. Cavalcante Siebert, M. L. Lupetti, E. Aizenberg, N. Beckers, A. Zgonnikov, H. Veluwenkamp, D. Abbink, E. Giaccardi, G.-J. Houben, C. M. Jonker, J. van den Hoven, D. Forster, and R. L. Lagendijk}\newline\vspace{0.2em}
\textit{AI and Ethics}\newline\vspace{0.2em}
\textbf{Year:} 2023\newline\vspace{0.2em}
\textbf{Publisher:} Springer Open Access}
&
Formalizes the philosophical concept of \textbf{Meaningful Human Control (\textbf{MHC})} into actionable software design properties: \textbf{Tracking} (system responsive to human moral reasons) and \textbf{Tracing} (system actions causally attributable to accountable human supervisors).
&
\textbf{Advantages:}
\begin{itemize}\itemsep0.05em\topsep0.1em
\item Translates abstract governance principles into rigorous software specifications.
\item Establishes tracking and tracing requirements for autonomous tool execution.
\item Directly guides the design of PrismSpace's Human-in-the-Loop (\textbf{HITL}) approval gates.
\end{itemize}
\vspace{-0.2em}
\textbf{Disadvantages:}
\begin{itemize}\itemsep0.05em\topsep0.1em
\item High human intervention frequency can degrade automated agent throughput.
\item Defining actionable reason boundaries requires careful domain calibration.
\item Assumes human supervisor maintains full comprehension of complex agent plans.
\end{itemize} \\
\hline
\end{tabular}
}
\end{frame}

\begin{frame}[t]{\textbf{CPN-HRL} (Cloud-Private-Network Hierarchical Reinforcement Learning): Hierarchical \textbf{DRL} for Priority-Aware Task Scheduling}
\vspace{-0.08cm}
\centering
\adjustbox{width=\textwidth, max totalheight=0.785\textheight, keepaspectratio}{%
\small
\renewcommand{\arraystretch}{1.12}
\setlength{\tabcolsep}{4pt}
\begin{tabular}{|p{3.1cm}|p{6.0cm}|p{6.0cm}|}
\hline
\textbf{Citation} & \textbf{Summary} & \textbf{Advantages \& Disadvantages} \\
\hline
\parbox[t]{\linewidth}{\raggedright\small
\textbf{Q. Yue, L. Tian, X. Feng, Z. Yuan, S. Wei, J. Xia, B. Chen, X. Song, and Y. Yao}\newline\vspace{0.2em}
\textit{Journal of King Saud University CIS}\newline\vspace{0.2em}
\textbf{Year:} 2026\newline\vspace{0.2em}
\textbf{Publisher:} Springer Open Access}
&
Presents \textbf{CPN-HRL}, a Hierarchical Reinforcement Learning framework for priority-aware task scheduling in Cloud-Edge-End networks. Uses a macro-policy to partition task \textbf{DAGs} by priority and latency budgets, while micro-policies schedule subtasks onto optimal compute nodes.
&
\textbf{Advantages:}
\begin{itemize}\itemsep0.05em\topsep0.1em
\item Hierarchical policy decomposition effectively handles complex task \textbf{DAGs}.
\item Guarantees priority-aware quality of service under tight latency deadlines.
\item Serves as primary reference for PrismSpace's \textbf{DAG} workflow scheduling engine.
\end{itemize}
\vspace{-0.2em}
\textbf{Disadvantages:}
\begin{itemize}\itemsep0.05em\topsep0.1em
\item Multi-tier policy architecture requires synchronous training coordination.
\item Edge node failure requires immediate micro-policy replanning.
\item Evaluated in simulated telecom cloud-edge topologies.
\end{itemize} \\
\hline
\end{tabular}
}
\end{frame}

\begin{frame}[t]{Automatic Discovery and Description of Human Planning Strategies}
\vspace{-0.25cm}
\centering
\adjustbox{width=\textwidth, max totalheight=0.785\textheight, keepaspectratio}{%
\small
\renewcommand{\arraystretch}{1.12}
\setlength{\tabcolsep}{4pt}
\begin{tabular}{|p{3.1cm}|p{6.0cm}|p{6.0cm}|}
\hline
\textbf{Citation} & \textbf{Summary} & \textbf{Advantages \& Disadvantages} \\
\hline
\parbox[t]{\linewidth}{\raggedright\small
\textbf{J. Skirzy\'{n}ski, Y. Jain, and F. Lieder}\newline\vspace{0.2em}
\textit{Behavior Research Methods}\newline\vspace{0.2em}
\textbf{Year:} 2023\newline\vspace{0.2em}
\textbf{Publisher:} Springer Open Access}
&
Develops an algorithmic framework for \textbf{automatically discovering and formalizing planning strategies} from human behavior traces. Models task planning as resource-rational search over Directed Graphical Models, producing interpretable decision trees that characterize human heuristic planning.
&
\textbf{Advantages:}
\begin{itemize}\itemsep0.05em\topsep0.1em
\item Extracts interpretable, human-like heuristic planning strategies automatically.
\item Quantifies trade-offs between cognitive search effort and decision quality.
\item Informs PrismSpace's hierarchical prompt decomposition and agent planning.
\end{itemize}
\vspace{-0.2em}
\textbf{Disadvantages:}
\begin{itemize}\itemsep0.05em\topsep0.1em
\item Search space for strategy discovery grows exponentially with problem depth.
\item Human behavioral experiments conducted in stylized cognitive testbeds.
\item Static decision trees struggle to capture fluid \textbf{LLM} reasoning trajectories.
\end{itemize} \\
\hline
\end{tabular}
}
\end{frame}

\begin{frame}[t]{Data-Intensive Workflow Scheduling Based on \textbf{DRL} in Multi-Clouds}
\vspace{-0.08cm}
\centering
\adjustbox{width=\textwidth, max totalheight=0.785\textheight, keepaspectratio}{%
\small
\renewcommand{\arraystretch}{1.12}
\setlength{\tabcolsep}{4pt}
\begin{tabular}{|p{3.1cm}|p{6.0cm}|p{6.0cm}|}
\hline
\textbf{Citation} & \textbf{Summary} & \textbf{Advantages \& Disadvantages} \\
\hline
\parbox[t]{\linewidth}{\raggedright\small
\textbf{S. Zhang, Z. Zhao, C. Liu, and S. Qin}\newline\vspace{0.2em}
\textit{Journal of Cloud Computing}\newline\vspace{0.2em}
\textbf{Year:} 2023\newline\vspace{0.2em}
\textbf{Publisher:} Springer Open Access}
&
Proposes a \textbf{Deep Q-Network (\textbf{DQN})} scheduler optimized for data-intensive workflows across multi-cloud environments. Explicitly models dataset placement and cross-datacenter transfer latency to minimize both workflow completion time and data egress billing.
&
\textbf{Advantages:}
\begin{itemize}\itemsep0.05em\topsep0.1em
\item Optimizes data locality to prevent massive cross-cloud egress overhead.
\item Adaptive reward function balances monetary cost and makespan constraints.
\item Directly relevant to caching and embedding retrieval across distributed agents.
\end{itemize}
\vspace{-0.2em}
\textbf{Disadvantages:}
\begin{itemize}\itemsep0.05em\topsep0.1em
\item Assumes deterministic file sizes and network bandwidth availability.
\item Q-network value overestimation requires double-\textbf{DQN} correction.
\item State representation size scales with the total number of cloud datacenters.
\end{itemize} \\
\hline
\end{tabular}
}
\end{frame}

\begin{frame}[t]{\textbf{RL}-Based Task Scheduling for Sustainable Federated Cloud Computing}
\vspace{-0.08cm}
\centering
\adjustbox{width=\textwidth, max totalheight=0.785\textheight, keepaspectratio}{%
\small
\renewcommand{\arraystretch}{1.12}
\setlength{\tabcolsep}{4pt}
\begin{tabular}{|p{3.1cm}|p{6.0cm}|p{6.0cm}|}
\hline
\textbf{Citation} & \textbf{Summary} & \textbf{Advantages \& Disadvantages} \\
\hline
\parbox[t]{\linewidth}{\raggedright\small
\textbf{Z. Wang, S. Chen, L. Bai, J. Gao, J. Tao, R. R. Bond, and M. D. Mulvenna}\newline\vspace{0.2em}
\textit{Journal of Cloud Computing}\newline\vspace{0.2em}
\textbf{Year:} 2023\newline\vspace{0.2em}
\textbf{Publisher:} Springer Open Access}
&
Develops an \textbf{Actor-Critic \textbf{RL} scheduling framework} for federated cloud datacenters that incorporates renewable energy availability into scheduling policies. Dynamically routes compute tasks to datacenters with green energy surpluses without violating user SLA deadlines.
&
\textbf{Advantages:}
\begin{itemize}\itemsep0.05em\topsep0.1em
\item Incorporates environmental sustainability and carbon metrics into scheduling.
\item Actor-Critic policy gradient handles continuous action spaces effectively.
\item Provides conceptual foundation for multi-attribute provider routing.
\end{itemize}
\vspace{-0.2em}
\textbf{Disadvantages:}
\begin{itemize}\itemsep0.05em\topsep0.1em
\item Requires real-time telemetry from green energy generation sources.
\item Trade-off between carbon footprint and network transmission latency.
\item Sensitive to reward scaling factors between SLA penalties and green rewards.
\end{itemize} \\
\hline
\end{tabular}
}
\end{frame}

\begin{frame}[t]{Hybrid Dolphin Swarm Sparrow Search Multi-Objective Workflow Scheduling}
\vspace{-0.25cm}
\centering
\adjustbox{width=\textwidth, max totalheight=0.785\textheight, keepaspectratio}{%
\small
\renewcommand{\arraystretch}{1.12}
\setlength{\tabcolsep}{4pt}
\begin{tabular}{|p{3.1cm}|p{6.0cm}|p{6.0cm}|}
\hline
\textbf{Citation} & \textbf{Summary} & \textbf{Advantages \& Disadvantages} \\
\hline
\parbox[t]{\linewidth}{\raggedright\small
\textbf{S. Subashree, M. Rajakumaran, G. Pushpa, and S. Senthilkumar}\newline\vspace{0.2em}
\textit{Journal of Cloud Computing}\newline\vspace{0.2em}
\textbf{Year:} 2025\newline\vspace{0.2em}
\textbf{Publisher:} Springer Open Access}
&
Proposes a hybrid bio-inspired metaheuristic combining \textbf{Dolphin Swarm Optimization (DSA)} and \textbf{Sparrow Search Algorithm (SSA)} for multi-objective cloud workflow scheduling. Optimizes execution makespan, virtual machine utilization, and monetary execution expenditure.
&
\textbf{Advantages:}
\begin{itemize}\itemsep0.05em\topsep0.1em
\item Combines global exploration of DSA with local exploitation of SSA.
\item Demonstrates superior Pareto hypervolume compared to standard \textbf{GA} and \textbf{PSO} (Particle Swarm Optimization).
\item Provides strong comparative baselines for workflow optimization algorithms.
\end{itemize}
\vspace{-0.2em}
\textbf{Disadvantages:}
\begin{itemize}\itemsep0.05em\topsep0.1em
\item Metaheuristic iteration runtime is too slow for sub-second reactive routing.
\item Multiple tuning parameters (predation rate, vocalization index) required.
\item Lacks theoretical convergence rate bounds for stochastic optimization.
\end{itemize} \\
\hline
\end{tabular}
}
\end{frame}

% Findings-transfer tables remain omitted from this presentation.

\begin{frame}{Research Gap}
\vspace{-0.05cm}
\begin{block}{What the surveyed methods already do}
\footnotesize
\textbf{MARL} works address \textbf{CTDE}, role grouping, and ad-hoc coordination. \textbf{ORPO} aligns \textbf{LLM} preferences without reference models, while \textbf{HIAT} leverages human interventions as auxiliary \textbf{RL} signals. \textbf{CPN-HRL} and \textbf{DAG} schedulers optimize cloud task allocation, and \textbf{MHC} formalizes tracking and tracing properties.
\end{block}
\vspace{0.08cm}
\begin{block}{Limitations observed in those papers}
\footnotesize
Coordination benchmarks focus on gaming (\textbf{SMAC}, Overcooked) rather than filesystem/\textbf{API} tooling. \textbf{ORPO} lacks grounding in \textbf{MCP} \textbf{ReAct} tool traces, while \textbf{HIAT} targets continuous robotic control instead of code execution. Cloud schedulers assume static \textbf{DAGs} or slow metaheuristics ill-suited for real-time \textbf{LLM} routing.
\end{block}
\vspace{0.08cm}
\begin{block}{Why PrismSpace is required}
\footnotesize
No reviewed system jointly grounds tool observations, aligns policies on agent trajectories, dynamically routes heterogeneous \textbf{LLMs}, and enforces human-in-the-loop safety gates within a unified developer workspace.
\end{block}
\end{frame}

\section{Project time line}

\begin{frame}{Gantt Chart}
\centering
\resizebox{1.0\textwidth}{!}{%
\begin{tikzpicture}[
  x=1.00cm, y=0.52cm,
  lab/.style={anchor=east, font=\scriptsize, inner sep=3pt, text=prismNavy}
]
\def\nm{9}
\def\rowh{1.0}
\def\bh{0.36}
\def\topY{11.5}

\foreach \i in {0,1,...,\nm} {
  \draw[gray!22] (\i,0.5) -- (\i,\topY);
}
\foreach \i/\m in {0.5/Jul'26,1.5/Aug'26,2.5/Sep'26,3.5/Oct'26,4.5/Nov'26,5.5/Dec'26,6.5/Jan'27,7.5/Feb'27,8.5/Mar'27} {
  \node[anchor=south, font=\scriptsize\bfseries, text=prismNavy] at (\i,\topY+0.08) {\m};
}
\draw[prismNavy, line width=0.7pt] (0,\topY) -- (\nm,\topY);
\draw[prismNavy, line width=0.6pt] (0,0.5) -- (\nm,0.5);

\foreach \k in {0,...,10} {
  \pgfmathsetmacro{\yc}{11.0-\k}
  \ifodd\k
    \fill[prismMint, opacity=0.45] (0,\yc-0.5) rectangle (\nm,\yc+0.5);
  \fi
}

\node[lab] at (0,11.0) {Project ideation \& problem identification};
\node[lab] at (0,10.0) {Literature survey \& research};
\node[lab] at (0,9.0) {Requirement analysis \& feasibility};
\node[lab] at (0,8.0) {System design (agents, \textbf{MCP}, \textbf{HITL})};
\node[lab] at (0,7.0) {Trajectory data collection \& pairing};
\node[lab] at (0,6.0) {Implementation: \textbf{MCP} \textbf{ReAct} swarm};
\node[lab] at (0,5.0) {Implementation: \textbf{ORPO} \& \textbf{ML} routing};
\node[lab] at (0,4.0) {Workspace integration \& \textbf{HITL} gates};
\node[lab] at (0,3.0) {Testing \& evaluation};
\node[lab] at (0,2.0) {Project documentation};
\node[lab] at (0,1.0) {Final review \& submission};

% Bars: same timeline, snapped to 0.5-month grid, equal height 0.36
\fill[prismBlue]     (0.0,11.0-0.18) rectangle (2.0,11.0+0.18);
\fill[prismTeal]     (1.0,10.0-0.18) rectangle (2.5,10.0+0.18);
\fill[prismGold]     (1.5,9.0-0.18)  rectangle (2.5,9.0+0.18);
\fill[prismBlue]     (2.0,8.0-0.18)  rectangle (3.5,8.0+0.18);
\fill[prismTeal]     (3.0,7.0-0.18)  rectangle (4.5,7.0+0.18);
\fill[prismGold]     (4.0,6.0-0.18)  rectangle (5.5,6.0+0.18);
\fill[prismBlue]     (5.0,5.0-0.18)  rectangle (6.5,5.0+0.18);
\fill[prismTeal]     (6.0,4.0-0.18)  rectangle (7.5,4.0+0.18);
\fill[prismGold]     (7.0,3.0-0.18)  rectangle (8.0,3.0+0.18);
\fill[prismBlue]     (7.5,2.0-0.18)  rectangle (9.0,2.0+0.18);
\fill[prismTeal]     (8.0,1.0-0.18)  rectangle (9.0,1.0+0.18);

\node[font=\scriptsize\bfseries, text=prismNavy] at (4.5,-0.15) {Timeline (Jul 2026 -- Mar 2027)};
\end{tikzpicture}%
}
\end{frame}

\section{References}

\begin{frame}{References}
\scriptsize
\setbeamertemplate{enumerate item}{\textbf{\color{prismBlue}[\insertenumlabel]}}
\begin{enumerate}\itemsep0.28em\leftskip=-0.55em
\item \textbf{Cornalba F, Disselkamp C, Scassola D, Helf C (2024)} Multi-objective reward generalization: improving performance of Deep Reinforcement Learning for applications in single-asset trading. \textit{Neural Computing and Applications}. DOI: 10.1007/s00521-023-09033-7
\item \textbf{Schneckenreither M, Moser G (2025)} Average reward adjusted discounted reinforcement learning. \textit{Neural Computing and Applications}. DOI: 10.1007/s00521-024-10620-5
\item \textbf{Aina K, Ha S (2026)} Deep reinforcement learning for multi-agent coordination. \textit{Artificial Life and Robotics}. DOI: 10.1007/s10015-025-01089-z
\item \textbf{Rother D, Pajarinen J, Peters J, Weisswange TH (2025)} Open-ended coordination for multi-agent systems using modular open policies. \textit{Autonomous Agents and Multi-Agent Systems}. DOI: 10.1007/s10458-025-09723-7
\item \textbf{Liu D, Ren F, Yan J, Su G, Gu W, Kato S (2026)} Improving scalability of multi-agent deep reinforcement learning with suboptimal human knowledge. \textit{Autonomous Agents and Multi-Agent Systems}. DOI: 10.1007/s10458-025-09729-1
\item \textbf{Marzari L, Liu C, Donti PL, Marchesini E (2026)} $\varepsilon$-retraining reinforcement learning algorithms. \textit{Autonomous Agents and Multi-Agent Systems}. DOI: 10.1007/s10458-026-09748-6
\item \textbf{Wong A, B\"ack T, Kononova AV, Plaat A (2023)} Deep multiagent reinforcement learning: challenges and directions. \textit{Artificial Intelligence Review}. DOI: 10.1007/s10462-022-10299-x
\item \textbf{Niu B, Luo B, Cui Y, Xu X, Zhao Y, Feng Y (2025)} HIAT: human-in-the-loop reinforcement learning with auxiliary task. \textit{Artificial Intelligence Review}. DOI: 10.1007/s10462-025-11260-4
\end{enumerate}
\end{frame}

\begin{frame}{References}
\scriptsize
\setbeamertemplate{enumerate item}{\textbf{\color{prismBlue}[\insertenumlabel]}}
\begin{enumerate}\setcounter{enumi}{8}\itemsep0.28em\leftskip=-0.55em
\item \textbf{Lee J (2025)} InstructPatentGPT: training patent language models to follow instructions with human feedback. \textit{Artificial Intelligence and Law}. DOI: 10.1007/s10506-024-09401-1
\item \textbf{Zhuang S, Qiao Z, Zuccon G (2023)} Reinforcement online learning to rank with unbiased reward shaping. \textit{Information Retrieval Journal}. DOI: 10.1007/s10791-022-09413-y
\item \textbf{Ghazal TM, Awouda AE, Hasan MK, Rahman AHA, Islam S, Saeed RA, Elshafie H, Iqbal A, Hussein AH (2025)} Knowledge Learning for Securing Workflow Scheduling Algorithm in Mobile Edge Computing. \textit{Mobile Networks and Applications}. DOI: 10.1007/s11036-025-02465-6
\item \textbf{Barredo P, Puente J (2025)} Energy-aware cooperative multi-fitness evolutionary algorithm for workflow scheduling in cloud computing. \textit{Natural Computing}. DOI: 10.1007/s11047-025-10023-y
\item \textbf{Heirati A, Fartash M, Khalily-Dermany M (2025)} Optimized Task Scheduling in Fog-Cloud Computing Using Hybrid Deep Learning and Metaheuristic Algorithms. \textit{Neural Processing Letters}. DOI: 10.1007/s11063-025-11819-w
\item \textbf{Abdi S, Ashjaei M, Mubeen S (2024)} Cost-aware workflow offloading in edge-cloud computing using a genetic algorithm. \textit{The Journal of Supercomputing}. DOI: 10.1007/s11227-024-06341-0
\item \textbf{Xia Y, Zhu J, Zhu L (2023)} Dynamic role discovery and assignment in multi-agent task decomposition. \textit{Complex \& Intelligent Systems}. DOI: 10.1007/s40747-023-01071-x
\item \textbf{Li J, Xing L, Zhong W, Cai Z, Hou F (2023)} Decision variable contribution based adaptive mechanism for evolutionary multi-objective cloud workflow scheduling. \textit{Complex \& Intelligent Systems}. DOI: 10.1007/s40747-023-01137-w
\end{enumerate}
\end{frame}

\begin{frame}{References}
\scriptsize
\setbeamertemplate{enumerate item}{\textbf{\color{prismBlue}[\insertenumlabel]}}
\begin{enumerate}\setcounter{enumi}{16}\itemsep0.22em\leftskip=-0.55em
\item \textbf{Yu X, Lin Y, Wang X, Han S, Lv K (2024)} GHQ: grouped hybrid Q-learning for cooperative heterogeneous multi-agent reinforcement learning. \textit{Complex \& Intelligent Systems}. DOI: 10.1007/s40747-024-01415-1
\item \textbf{Xu C, Wang J, Zhu X, Yue Y, Zhou W, Liang Z, Wojtczak D (2024)} Decentralized multi-agent cooperation via adaptive partner modeling. \textit{Complex \& Intelligent Systems}. DOI: 10.1007/s40747-024-01421-3
\item \textbf{Vyhmeister E, Casta\~n\'e GG, \"Ostberg PO, Thevenin S (2023)} A responsible AI framework: pipeline contextualisation. \textit{AI and Ethics}. DOI: 10.1007/s43681-022-00154-8
\item \textbf{Cavalcante Siebert L, Lupetti ML, Aizenberg E, Beckers N, Zgonnikov A, Veluwenkamp H, Abbink D, Giaccardi E, Houben GJ, Jonker CM, van den Hoven J, Forster D, Lagendijk RL (2023)} Meaningful human control: actionable properties for AI system development. \textit{AI and Ethics}. DOI: 10.1007/s43681-022-00167-3
\item \textbf{Yue Q, Tian L, Feng X, Yuan Z, Wei S, Xia J, Chen B, Song X, Yao Y (2026)} CPN-HRL: a hierarchical deep reinforcement learning approach for priority-aware task scheduling in CPN enabled by cloud-edge-end environments. \textit{Journal of King Saud University Computer and Information Sciences}. DOI: 10.1007/s44443-026-00690-x
\item \textbf{Skirzy\'{n}ski J, Jain Y, Lieder F (2023)} Automatic discovery and description of human planning strategies. \textit{Behavior Research Methods}. DOI: 10.3758/s13428-023-02062-z
\item \textbf{Zhang S, Zhao Z, Liu C, Qin S (2023)} Data-intensive workflow scheduling strategy based on deep reinforcement learning in multi-clouds. \textit{Journal of Cloud Computing}. DOI: 10.1186/s13677-023-00504-9
\item \textbf{Wang Z, Chen S, Bai L, Gao J, Tao J, Bond RR, Mulvenna MD (2023)} Reinforcement learning based task scheduling for environmentally sustainable federated cloud computing. \textit{Journal of Cloud Computing}. DOI: 10.1186/s13677-023-00553-0
\item \textbf{Subashree S, Rajakumaran M, Pushpa G, Senthilkumar S (2025)} Hybrid dolphin swarm sparrow search optimization based multi-objective cloud workflow scheduling. \textit{Journal of Cloud Computing}. DOI: 10.1186/s13677-025-00828-8
\end{enumerate}
\end{frame}

\begin{frame}
    \centering
    \vfill
    {\usebeamerfont{title}\color{prismNavy}\Huge \textbf{Thank You}}
    
    \vspace{0.8cm}
    {\large \textbf{Questions \& Discussion}}
    
    \vspace{0.8cm}
    {\small \color{prismBlue} PrismSpace: An Intelligent Developer Workspace for Multi-Agent Orchestration}
    \vfill
\end{frame}
